\chapter{Memoria torsional y condiciones de rebote en Einstein--Cartan}
\label{ch:rebote-ec}

\section{Transporte de la ley de Perron a la densidad espectral de espín}

La realizaci\'on cosmol\'ogica de Perron construye dos secuencias sobre una
misma ruta cerrada:

\begin{equation}
 \frac{a_n}{a_0}=\varphi^{n/3},
 \qquad
 M_n=\varphi^{-2n}.
 \label{eq:rebote-perron-pair}
\end{equation}

La eliminaci\'on del nivel \(n\) produce la identidad exacta

\begin{equation}
 \boxed{M_n=\left(\frac{a_n}{a_0}\right)^{-6}.}
 \label{eq:rebote-memory-six}
\end{equation}

El exponente posee una genealog\'ia determinada: la dimensi\'on espacial
aporta el factor \(3\) y el car\'acter cuadr\'atico de la memoria estable
aporta el factor \(2\).  La potencia \(6=2\cdot3\) expresa as\'i la diluci\'on
de una densidad cuadr\'atica asociada a una carga extensiva de esp\'in.

\begin{definition}[Transductor torsional]
Fijada una amplitud positiva \(s_*^2\) en una secci\'on inicial, el transductor
torsional es
\begin{equation}
 \mathcal T_{\rm EC}:M_n\longmapsto
 s_n^2=s_*^2M_n
 =s_*^2\left(\frac{a_n}{a_0}\right)^{-6}.
 \label{eq:rebote-transductor}
\end{equation}
Su dominio es el modo estable de Perron y su codominio es el sector
cuadr\'atico de esp\'in de una realizaci\'on de Einstein--Cartan.
\end{definition}

La ecuaci\'on \cref{eq:rebote-transductor} conserva el exponente y la
positividad.  La amplitud \(s_*^2\), el signo con que la contorsi\'on entra en
la ecuaci\'on gravitatoria y la normalizaci\'on dimensional pertenecen al
realizador f\'isico.  Esta separaci\'on sit\'ua exactamente qu\'e aporta el
tronco HMT y qu\'e debe aportar la din\'amica de Cartan.

\section{Ecuaciones homog\'eneas y convenci\'on de densidad}

Para exponer el mecanismo estructural adoptamos temporalmente unidades
\(c=1\) y una geometr\'ia FRW espacialmente plana.  Sea

\begin{equation}
 \rho_m(a)=\rho_0a^{-3(1+w)},
 \qquad
 \rho_s(a)=s_0a^{-6},
 \label{eq:rebote-densities}
\end{equation}

con \(\rho_0>0\), \(s_0>0\) y \(w\) constante.  La contribuci\'on de
contorsi\'on entra con signo repulsivo en la ecuaci\'on efectiva:

\begin{equation}
 H^2=\frac{8\pi G}{3}\bigl(\rho_m-\rho_s\bigr).
 \label{eq:rebote-friedmann}
\end{equation}

Si se emplean densidades de energ\'ia, el lado derecho se multiplica por
\(c^{-2}\); si se emplean densidades de masa, conserva la forma
\cref{eq:rebote-friedmann}.  El documento mantiene una sola convenci\'on en
cada c\'alculo y restaura la carta dimensional al final.

La ecuaci\'on de Raychaudhuri correspondiente es

\begin{equation}
 \dot H=-4\pi G\bigl[(1+w)\rho_m-2\rho_s\bigr].
 \label{eq:rebote-raychaudhuri}
\end{equation}

La segunda contribuci\'on se comporta como un fluido r\'igido de ecuaci\'on de
estado \(w_s=1\) y densidad efectiva negativa.  Esta forma es la publicaci\'on
homog\'enea de la repulsi\'on torsional a gran densidad.

\section{Existencia y unicidad del radio de rebote}

\begin{theorem}[Rebote regular para \(w<1\)]
Sup\'onganse \(\rho_0>0\), \(s_0>0\) y \(w<1\).  Entonces existe un \'unico
radio positivo \(a_b\) tal que \(H(a_b)=0\), dado por
\begin{equation}
 \boxed{
 a_b^{3(1-w)}=\frac{s_0}{\rho_0}.}
 \label{eq:rebote-radius}
\end{equation}
En ese punto
\begin{equation}
 \boxed{
 \dot H_b=4\pi G(1-w)\rho_b>0,}
 \qquad
 \rho_b:=\rho_m(a_b)=\rho_s(a_b),
 \label{eq:rebote-positive-derivative}
\end{equation}
de modo que la rama contractiva se prolonga en una rama expansiva regular.
\end{theorem}

\begin{proof}
La condici\'on \(H=0\) equivale a
\(\rho_0a^{-3(1+w)}=s_0a^{-6}\).  Multiplicar por \(a^6\) produce
\cref{eq:rebote-radius}.  El exponente \(3(1-w)\) es positivo; por tanto la
potencia define una biyecci\'on de \(\RR_{>0}\) en s\'i misma y el radio es
\'unico.  Sustituyendo \(\rho_m=\rho_s=\rho_b\) en
\cref{eq:rebote-raychaudhuri},
\[
 \dot H_b=-4\pi G[(1+w)-2]\rho_b
 =4\pi G(1-w)\rho_b>0.
\]
La positividad convierte el cero simple de \(H\) en un m\'inimo de \(a(t)\).
\end{proof}

La expansi\'on local alrededor de \(t_b\) hace visible la regularidad:

\begin{equation}
 a(t)=a_b\left[
 1+2\pi G(1-w)\rho_b(t-t_b)^2
 +O\!\left((t-t_b)^3\right)
 \right].
 \label{eq:rebote-local-expansion}
\end{equation}

La ley \(a^{-6}\) aporta la potencia que domina al sector material cuando
\(w<1\) y \(a\) disminuye.  El signo torsional y la amplitud positiva hacen
que esa dominancia cierre la contracci\'on antes de alcanzar \(a=0\).

\section{Casos límite de la dinámica torsional y obstrucciones a la regularidad del rebote}

El valor \(w=1\) produce el mismo exponente \(a^{-6}\) en ambos t\'erminos.
La diferencia conserva entonces un signo fijo determinado por
\(\rho_0-s_0\); la igualdad puntual deja de seleccionar un radio aislado.
Para \(w>1\), la densidad material crece m\'as deprisa que el t\'ermino
torsional durante la contracci\'on y el mecanismo anterior no garantiza un
rebote.  Estos casos proporcionan controles negativos internos: el cierre
depende de la desigualdad de exponentes y no de una elecci\'on retrospectiva
del umbral.

La curvatura espacial \(k\), una constante cosmol\'ogica y las anisotrop\'ias
de Bianchi a\~naden t\'erminos a \cref{eq:rebote-friedmann}. Su efecto local
se controla sin presuponer la posici\'on del nuevo cero. Definamos
\[
 F_0(a):=\frac{8\pi G}{3}\bigl(\rho_m(a)-\rho_s(a)\bigr).
\]
La ecuaci\'on perturbada adopta la forma
\(H^2=F_R(a):=F_0(a)+R(a)\).

\begin{theorem}[Persistencia \(C^1\) del rebote transversal]
\label{thm:rebote-c1-persistence}
Sea \(a_b\) el radio de \cref{eq:rebote-radius}. Existen
\(0<a_-<a_b<a_+\) tales que, para
\(I=[a_-,a_+]\),
\[
 F_0(a_-)<0<F_0(a_+),
 \qquad
 m_I:=\inf_{a\in I}F_0'(a)>0.
\]
Si \(R\in C^1(I)\) satisface
\begin{equation}
 \|R\|_{C^0(I)}
 <\delta_I
 :=\min\{-F_0(a_-),F_0(a_+)\},
 \qquad
 \|R'\|_{C^0(I)}<m_I,
 \label{eq:rebote-c1-bounds}
\end{equation}
entonces existe un único \(a_R\in(a_-,a_+)\) con \(F_R(a_R)=0\), y
ese cero es transversal:
\begin{equation}
 F_R'(a_R)>0,
 \qquad
 \dot H_R(a_R)=\frac{a_R}{2}F_R'(a_R)>0.
 \label{eq:rebote-perturbed-transversality}
\end{equation}
Para una familia \(R(a,\eta)\) de clase \(C^1\), el teorema de la funci\'on
impl\'icita proporciona una rama local \(a_R(\eta)\) de clase \(C^1\).
\end{theorem}

\begin{proof}
La identidad \cref{eq:rebote-positive-derivative} equivale a
\(F_0'(a_b)>0\). Por continuidad puede escogerse \(I\) con
\(m_I>0\) y con el cambio de signo indicado. La primera cota de
\cref{eq:rebote-c1-bounds} conserva
\(F_R(a_-)<0<F_R(a_+)\); el teorema del valor intermedio produce un cero.
La segunda cota da
\(F_R'(a)\geq m_I-\|R'\|_{C^0(I)}>0\), de modo que \(F_R\) es estrictamente
creciente y el cero es único. En cada rama con \(H\neq0\), derivar
\(H^2=F_R(a)\) y usar \(\dot a=aH\) da
\(\dot H=(a/2)F_R'(a)\); la continuidad extiende la igualdad a \(a_R\).
La no degeneración \(F_R'(a_R)>0\) es precisamente la hipótesis del
teorema de la función implícita.
\end{proof}

\section{Amplitud torsional desde la pantalla de Cartan}

El exponente de la memoria queda cerrado por
\cref{eq:rebote-memory-six}.  La amplitud puede buscarse prospectivamente en
la respuesta de pantalla.  Sea la identidad exacta

\begin{equation}
 (\beta_m^{\square})^*\Lambda_m^{\square}\beta_m^{\square}
 =\frac{56}{81}I_{Q_{\square}}.
 \label{eq:rebote-screen-energy}
\end{equation}

Una realizaci\'on gravitatoria completa debe construir una corriente de esp\'in

\begin{equation}
 \mathcal J_m
 =\mathcal J(\beta_m^{\square},U_m,c_s,c_g)
 \label{eq:rebote-spin-current}
\end{equation}

y un funcional positivo

\begin{equation}
 s_{0,m}=\mathfrak S(\mathcal J_m,\Lambda_m^{\square})
 \label{eq:rebote-amplitude-functional}
\end{equation}

compatible con refinamiento.  La normalizaci\'on non\'adica pertinente es la
esperanza tracial, porque conserva la densidad de
\cref{eq:rebote-screen-energy}; la suma bruta cuenta nueve copias y pertenece
a una magnitud extensiva distinta.

\begin{proposition}[Obligaciones del selector de amplitud]
Un selector admisible para \(s_0\) debe satisfacer:
\begin{enumerate}[label=\roman*)]
 \item positividad y dimensiones de densidad de esp\'in al cuadrado;
 \item invariancia gauge bajo la reducci\'on c\'ubica transportada;
 \item compatibilidad con la inversi\'on de arista y el carry;
 \item naturalidad bajo el refinamiento \(9\)-ario;
 \item independencia de cualquier radio de rebote usado como valor objetivo.
\end{enumerate}
\end{proposition}

Estas condiciones convierten la amplitud en una obligaci\'on algebraica
localizada.  La igualdad energ\'etica de pantalla suministra la forma
cuadr\'atica; el realizador de Cartan debe determinar c\'omo esa forma se
contrae con la corriente f\'isica.

\section{Torsi\'on, orientaci\'on y memoria de transporte}

La m\'etrica publica distancias y conos causales.  La conexi\'on de Cartan
publica adem\'as la orientaci\'on del transporte.  En HMT, ruta, carry y
estado profundo constituyen ya datos de la genealog\'ia.  La correspondencia

\begin{equation}
 \text{memoria estable de ruta}
 \longrightarrow
 \text{corriente de esp\'in}
 \longrightarrow
 \text{contorsi\'on}
 \longrightarrow
 \rho_s\propto a^{-6}
 \label{eq:rebote-genealogy}
\end{equation}

ofrece una lectura f\'isica de esa informaci\'on orientada.  La curvatura
describe la respuesta de la geometr\'ia a la energ\'ia; la torsi\'on conserva
la historia de orientaci\'on que acompa\~na al transporte.  La ley de sexto
orden es el punto en que la memoria geneal\'ogica y la densidad cuadr\'atica de
esp\'in poseen exactamente la misma covariancia de escala.

\section{Principio de Mach y autoinercia cosmográfica}

El estado local se eval\'ua dentro de una ruta compatible con la frontera
global.  La respuesta inercial pertenece, por ello, al estabilizador
holon\'omico de la totalidad.  En el sector homog\'eneo, la escala \(a_n\), la
memoria \(M_n\) y la corriente torsional son tres lecturas del mismo estado
geneal\'ogico enriquecido:

\begin{equation}
 X_n\longmapsto
 \bigl(a_n,M_n,\mathcal J_n\bigr).
 \label{eq:rebote-mach-triple}
\end{equation}

La interpretaci\'on machiana adquiere as\'i un contenido operatorio: la
respuesta local depende de la clase global de cierre y el observador forma
parte de la secci\'on prolongada.  La autoinercia es la conservaci\'on del
estado relativo bajo la holonom\'ia que transporta conjuntamente geometr\'ia,
observador y memoria.

\section{Identidades exactas de la dinámica torsional y condiciones de la interfaz Einstein–Cartan}

\begin{proofstatusbox}
La ley \(M_n=(a_n/a_0)^{-6}\) es un teorema interno exacto.  El teorema de
rebote \cref{eq:rebote-radius,eq:rebote-positive-derivative} es una
consecuencia exacta de las ecuaciones homog\'eneas de Einstein--Cartan
\cite{Hehl1976} con
\(\rho_0,s_0>0\) y \(w<1\).  El transductor
\cref{eq:rebote-transductor} conserva exactamente la ley de escala.  La
determinaci\'on de \(s_0\) desde la corriente de pantalla, la extensi\'on a
anisotrop\'ias y la derivaci\'on completa de la acci\'on Einstein--Cartan son
interfaces f\'isicas delimitadas.
\end{proofstatusbox}

\begin{falsifierbox}
El bloque interno se refuta si el modo de Perron incumple
\cref{eq:rebote-memory-six}.  La realizaci\'on de Einstein--Cartan se refuta
si la corriente inducida produce otra potencia, una amplitud no positiva o
un signo que elimina \cref{eq:rebote-positive-derivative}.  El rebote
homog\'eneo falla para los datos declarados si no existe el radio
\cref{eq:rebote-radius}, si \(H(a_b)\ne0\) o si \(\dot H_b\le0\).  La
promoci\'on completa queda rechazada si el funcional de amplitud consulta
retrospectivamente el radio deseado o pierde gauge, carry o refinamiento.
\end{falsifierbox}
